\section{Egonet Bags Among Node Embeddings}
\label{sec:baselines}

The experiments so far compare egonet bags against two references only:
the permutation floor and the node's own feature profile. This section
places them among the classic node-embedding methods --- and uses the
instrument of Section~\ref{sec:structural-share} to ask of every method
the question that section asks of ours: how much of its performance is
structure, and how much is position? Node embeddings fall into two
families that differ in exactly this respect. \emph{Proximity} methods ---
DeepWalk \citep{perozzi2014deepwalk}, node2vec \citep{grover2016node2vec},
and their matrix-factorization relatives \citep{qiu2018netmf} --- place
nodes that co-occur on short random walks close together, so the
embedding encodes \emph{where} a node sits in the graph; under a
spatially smooth label field this is the locality-leakage pathway of
Section~\ref{sec:synthetic} turned into a representation. \emph{Structural}
(role) methods --- struc2vec \citep{ribeiro2017struc2vec}, Role2Vec
\citep{ahmed2018role2vec}, GraphWave \citep{donnat2018graphwave}, and the
recursive-feature lineage of ReFeX/RolX
\citep{henderson2011refex,henderson2012rolx,rossi2015roles} --- place nodes
with similar \emph{surroundings} close together regardless of distance,
and are position-free by construction. Egonet bags belong, in intent, to
the second family. The comparison therefore doubles as a test of the
mechanism account: the same region-holdout evaluation that separated
structure from position for our bags should separate the two embedding
families from each other.

\paragraph{Methods and protocol.}
%% (drafted before results; numbers and findings filled from
%% outputs/stats_*.csv via make_tables.py)
We fit six embeddings on every dataset. DeepWalk and node2vec (two
biased settings: breadth-first $p{=}0.25,\,q{=}4$ and depth-first
$p{=}4,\,q{=}0.25$) use second-order walks with the exact transition
probabilities of \citet{grover2016node2vec} and a skip-gram
\citep{mikolov2013word2vec} trained with gensim \citep{rehurek2010gensim}.
struc2vec, Role2Vec and GraphWave have no maintained implementation we
were willing to depend on, so we implemented each from its paper:
struc2vec with the authors' optimizations (compressed degree sequences
with exact dynamic time warping, $2\lfloor\log_2 n\rfloor{+}1$
similar-degree candidates per node, hierarchy capped at six layers;
$q{=}0.3$ as in the reference code); Role2Vec in its original form
(structural features --- degree, triangles, open wedges, core number ---
log-binned into node types, walks rewritten as type sequences, one vector
per type); GraphWave following the reference implementation (heat
wavelets of the normalized Laplacian at the reference's automatic scales,
Chebyshev-approximated, empirical characteristic function at 25 sample
points, 100 features). Every learned embedding has width $D{=}16$ ---
wider than the bags' $D{=}11$, so dimension cannot favor our method ---
and all walk methods use 10 walks of length 80 with window 10. Each
embedding is fit unsupervised on the full graph, exactly as the bags and
the node features are, and is then evaluated on the identical center sets
and the identical ten stratified splits of
Section~\ref{sec:experiments}, with a permutation floor computed per
method; all comparisons are therefore paired. Predictions were registered
before evaluation: proximity embeddings should be strongest on the
homophilous, position-rich graphs and collapse under region holdout;
structural embeddings should be strongest where the label is a role
(Air Traffic is struc2vec's own benchmark) and survive separation.

%% RESULTS — to be written from the CSVs
%% \input{tables/baselines_shared}  figure baselines_overview
%% \input{tables/baselines_holdout} figure baselines_separation
